%-------------------------------------------------------------------------------
%	ABSTRACT PAGE
%-------------------------------------------------------------------------------

\addchaptertocentry{Хураангуй} % Хураангуйг гарчигт нэмэх

\begin{center}
{\scshape\Large \univname\par} % Их сургуулийн нэр
{\scshape\large \facname\par}\vspace{0.5cm} % Их сургуулийн нэр
{\huge\textbf{{Хураангуй}} \par}
\bigskip
{\Large{\ttitle} \par} % Тезисийн нэр
\bigskip

{\normalsize \shortname \par} % Зохиогчийн нэр
\addressname
\end{center}

\textit{\textbf{Түлхүүр үгс: \keywordnames}}
\bigskip

Өнөөгийн мэдээллийн эрин үед богино бичвэр мессеж нь хүмүүсийн өдөр тутмын харилцааны гол хэрэгсэл болоод байна. Гэсэн хэдий ч ирж буй мессежийн хэмжээ нэмэгдэхийн хэрээр хэрэглэгчийн ирсэн мессежийн жагсаалт замбараагүй болж, чухал мэдээллийг олоход хүндрэлтэй болох асуудал тулгарсаар байна. Энэхүү судалгааны ажлын зорилго нь уг асуудлыг шийдвэрлэхийн тулд мессежийг агуулгын дагуу автоматаар ангилах болон хавтаслах систем боловсруулах явдал юм.


Судалгааны хүрээнд нээлттэй эхийн Fossify Messages аппликейшнд суурилан Android платформ дээр шинэ функциональ боломжуудыг хөгжүүлсэн. Хөгжүүлэлтэд MVVM архитектурын загвар, Room өгөгдлийн сангийн өргөтгөл, Jetpack Compose-д суурилсан түлхүүр үг удирдлагын систем болон өргөн дэлгэцтэй төхөөрөмжид зориулсан дасан зохицох харагдацыг хэрэгжүүллээ. Мөн мессежийг автоматаар таних, ангилах, хавтасанд байршуулах дэс дараалсан алгоритм боловсруулж, Room өгөгдлийн сангийн шилжилтийн тестээр баталгаажуулсан.


Хийгдсэн ажлын үр дүнд хэрэглэгч мессежээ гараар эрж хайхгүйгээр автоматаар ангилагдсан байдлаар харах, хавтасаар шүүх боломжтой болсон бөгөөд орчин үеийн мессеж аппликейшнуудын давуу талуудыг нэгтгэсэн шийдлийг нээлттэй эхийн орчинд хэрэгжүүлж чадсан юм.


